\scanpage{30}
7. \uline{Stratifications} et \uline{voisinages côniques} : le pas
de récurrence.

$(X_i)_{i\in I}$, famille croissante de parties de $X$,
satisfaisant
\[
(14)\qquad\text{a) }\forall i,j\in I,\ \text{on a}\qquad
X_i\cap X_j=\bigcup_{\alpha\le i,j}X_\alpha.
\]
\inserted{$X=\bigcup_{i\in I}X_i$}

\margin{Cette condition se généralise :
aux intersections quelconques
$\displaystyle\bigcap_{i\in J}X_i=
\bigcup_{\substack{k\in I\\k\text{ minorant }J}}X_k$.}

Transversalité : un syst. de parties-bord \inserted{$B_\alpha$} de
$X$, munies de germes de cylindres associés
(mutuellement transversaux entre eux) : si on
représente les germes par des cylindres convenables
$E_{B_\alpha}$, \unsure{imposant}\inserted{\unsure{exprimant}} que les $X_i$ soient transversaux
aux $E_{B_\alpha}$ --- cela implique qu'on a
\[
E_J\cap X_i\simeq (X_{iJ}\times I^J
\qquad
\text{(induit par }E_J\simeq B_J\times I^J,
\]
avec $X_{iJ}=X_i\cap B_J$).
\deleted{\unsure{Bien sûr les}}
\margin{pour $J$ \deleted{\unsure{ensemble}} \unsure{ensemble} d'indices $\alpha$}

\deleted{\unsure{Soit}}
Soit $I_0\subset I$ une partie de $I_0$ formée
d'éléments minimaux, \deleted{et \unsure{soit}} et soit
\[
X_{I_0}=\bigcup_{i\in I_0}X_i=\coprod_{i\in I_0}X_i.
\]
\deleted{\unsure{Alors}}
La donnée d'un germe \deleted{de structure}
d'un voisinage cônique \inserted{strict} de $X_{I_0}$ équivaut à
celle d'un syst de voisinages côniques
$T_i$ ($i\in I_0$) mutuellement disjoints dans $X$. ---
celle d'un germe de voisinage cônique de $X$,
à celle d'une famille de germes de vois.
côniques des $X_i$ dans $X$.
\margin{29}
